% Generated draft: requires mathematical and source review.
\documentclass[11pt]{article}
\usepackage[margin=25mm]{geometry}
\usepackage{fontspec}
\setmainfont{Latin Modern Roman}
\usepackage{amsmath,amssymb,mathtools,unicode-math}
\setmathfont{Latin Modern Math}
\usepackage{enumitem,needspace,xurl,hyperref}
\setlength{\parindent}{0pt}
\setlength{\parskip}{5pt}
\setlength{\emergencystretch}{3em}
\newcommand{\N}{\mathbb{N}}
\newcommand{\Z}{\mathbb{Z}}
\newcommand{\Q}{\mathbb{Q}}
\newcommand{\R}{\mathbb{R}}
\newcommand{\C}{\mathbb{C}}
\newcommand{\F}{\mathbb{F}}
\newcommand{\A}{\mathbb{A}}
\newcommand{\PP}{\mathbb{P}}
\newcommand{\E}{\mathbb{E}}
\newcommand{\eps}{\varepsilon}
\newcommand{\dd}{\,\mathrm{d}}
\begin{document}
\section*{UnsolvedMath ID: 700935}
\textbf{Source problem code:} catalogue entry
\subsection{Definitions and mathematical statement}
Let \(K\) be a field of characteristic zero, let \(G\) be a group, and let \(V\ne\{0\}\) be a right module over the group algebra \(K[G]\). No finite-dimensionality assumption is imposed. For \(g\in G\), the expression \(v(g-1)^n\) denotes the action of the indicated group-algebra element on \(v\).

The representation is a \emph{nil-representation} if
\[
\forall v\in V\;\forall g\in G\;\exists n=n(v,g)\in\Z_{\ge1}
\quad v(g-1)^n=0.
\]
It is \emph{irreducible} if its only \(K\)-linear subspaces invariant under every element of \(G\) are \(\{0\}\) and \(V\). It is \emph{trivial} if \(vg=v\) for every \(v\in V\) and \(g\in G\).

\emph{Question (S. Vovsi).} For every field \(K\) of characteristic zero, every group \(G\), and every nonzero irreducible nil-representation \(V\) over \(K\), must the representation be trivial? Equivalently, must
\[
vg=v\qquad\text{for all }v\in V\text{ and }g\in G?
\]
The supplied source additionally states that the analogous assertion is false in prime characteristic; it supplies no counterexample or further qualification for that statement.
\subsection{Short English statement}
If each group element acts locally unipotently on a nonzero irreducible representation over a field of characteristic zero, must every group element act as the identity? The exponent may depend on both the vector and the group element, and the representation need not be finite-dimensional.
\subsection{Sources}
\begin{itemize}
\item Local UnsolvedMath registry, permanent ID \(700935\), source problem code \texttt{catalogue entry}, supplied complete source record; upstream revision \texttt{d371ee917ac9647fbb06cd6fc529deec6f3dda82}.
\item \emph{The Kourovka Notebook (Unsolved Problems in Group Theory)}, Problem 13.17, proposed by S. Vovsi, 13th issue (1995); cited in the supplied record through the 21st-issue edition (2026), \texttt{arXiv:1401.0300v48}. The collection metadata names E. I. Khukhro and V. D. Mazurov as editors. These bibliographic details are supplied attribution and have not been independently verified.
\end{itemize}
\end{document}
